% Pulsar Notes on Patent Claims -- included by pulsar.tex.
% Mirrors docs/patent-notes-draft.md (now deleted; this file is the
% submission-package patent-claims appendix).

\paragraph{Required by NIST MPTC IR 8214C \S5} as a package
deliverable. This document is collected over the development period
and finalized for the submission.

\subsection{Status}

\textbf{Draft.} No patent claims have been formally identified. This
appendix tracks claims as contributors disclose them.

\subsection{Inherited from upstream}

\paragraph{Corona (\texttt{luxfi/corona}).}
Pulsar reuses Corona's protocol structure without importing its code.
Corona is distributed under the Lux Ecosystem License v1.2, which
permits research use and reserves commercial use to a separate licence
from Lux Industries Inc. It is not an Apache-2.0 grant and carries no
Apache~\S3 patent grant.

Corona's \texttt{PATENTS.md} asserts no Lux patents on the published
Boschini et al.\ algorithm itself, which it implements unchanged in its
mathematics.

\paragraph{Ringtail \cite{boschini2024ringtail}.}
Both Corona and Pulsar descend from Ringtail, the two-round threshold
construction of Cecilia Boschini, Darya Kaviani, Russell W. F. Lai,
Giulio Malavolta, Akira Takahashi and Mehdi Tibouchi (ePrint 2024/1113,
IEEE S\&P 2025). Corona implements it over Ring-LWE; Pulsar retargets
the same structure to Module-LWE for FIPS~204. Ringtail is the authors'
name for their scheme; Corona and Pulsar are ours for our
implementations.

\paragraph{ML-DSA / FIPS 204.}
ML-DSA is a NIST-standardized algorithm. The Dilithium team's IP
disclosures during NIST PQC standardization are public on the NIST
PQC website. Pulsar targets output interchangeability with
FIPS~204 but does not assert any rights over FIPS~204 itself.

\subsection{Contributor disclosures}

Per \texttt{CONTRIBUTING.md}, every commit's \texttt{Signed-off-by}
line is taken as either:

\begin{enumerate}[leftmargin=*, itemsep=0pt]
\item An explicit assertion that the contributor has no relevant
patent claims, OR
\item A reference to a disclosed claim listed below.
\end{enumerate}

\begin{center}
\begin{tabular}{llll}
\toprule
date & contributor & type & disclosure \\
\midrule
2026-05-10 & Lux Industries Inc & initial repo & no claims asserted \\
\bottomrule
\end{tabular}
\end{center}

(Updated as contributions land.)

\subsection{Third-party patents we may need to navigate}

Live concerns to investigate before submission:

\begin{itemize}[leftmargin=*, itemsep=0pt]
\item Lattice-estimator usage (we use the open-source estimator;
license pass-through to MIT).
\item Gaussian sampler implementation (FIPS~204 references CDT;
alternate rejection samplers may have IP).
\item Polynomial multiplication via NTT (standard, no known IP).
\item Pedersen commitments (1991 paper, expired patents).
\end{itemize}

The submission's Notes on Patent Claims will explicitly call out
each of these.

\subsection{NIST MPTC submission requirements}

NIST IR 8214C \S5 (Notes on Patent Claims) requires:

\begin{quote}
Submitter(s) shall include a statement of any patent claims known
to them that may be relevant to the submitted package, regardless
of who owns the claim. Statements of ``no known claims'' are
acceptable.
\end{quote}

Our final statement, once the spec freezes:

\begin{quote}
\emph{To be drafted at spec freeze. Will incorporate every
disclosure in the contributor table above, every inherited claim
from Pulsar / Corona / ML-DSA, and the third-party investigation
outcome.}
\end{quote}

\subsection{Process}

\begin{enumerate}[leftmargin=*, itemsep=0pt]
\item Every PR review checks for new IP-relevant code.
\item Significant cryptographic additions (new sampling routine, new
commitment shape, new aggregation step) trigger an explicit IP
review with one of the sponsor's IP counsel.
\item Final statement assembled at freeze, included in
\texttt{spec/pulsar.pdf} Appendix~A.
\end{enumerate}
